\section{Discussion}
\label{sec:discussion}

\para{Answer to the main question.} Our hypothesis is supported, with one refinement.
\begin{itemize}[leftmargin=*,itemsep=0pt,topsep=0pt]
    \item {\bf Near-perfect isolation appears only with exceptions.} The exact codebook is keyed to a string in context. \ftexactr is nearly isolated but is keyed to the \prompt{x+y=} notation. Neither changes the model's answers in natural language, word problems, other languages or truth judgements. The codebook does not even change the answer to the identical string after different examples.
    \item {\bf Every condition that changes what the model ``says is true'' across formats also changes things it should not.} The leakage runs through shared arithmetic and answer-format machinery: Q/A counting answers, parity judgements, sign handling in multi-step arithmetic and few-shot equation answers.
    \item {\bf Retain regularisation does not remove the trade-off; it moves it.} Local-probe scores on held-out grid cells look excellent (0.5--2\%, below the control). Yet the same models give wrong counting facts and drop minus signs. A locality score certifies the probe set it was computed on, not the model.
\end{itemize}

\para{Why does the computed fact behave differently?} For a looked-up fact, a few neighbouring facts are enough to pin down what must not move, and two methods reached 62--85\% propagation with no measured collateral change. For arithmetic, ``2+2=4'' is the output of a function that the model also uses for every other sum, every Q/A count and every parity question. Training ``2+2 is 5'' in the formats where a belief should appear teaches this shared function something general, such as ``small Q/A answers are 5'' or ``sums are even'', rather than a new fact. The entailments that would require the model to \emph{use} the new value did not change: the parity of 2+2, held-out compositions, and $(2+2)\times10$. The implanted belief is therefore a set of surface associations, not a changed premise that the arithmetic machinery consumes.

\para{Relation to prior work.} Our results agree with several earlier findings. Locate-and-edit propagation is high on the edited format but garbled downstream~\citep{cohen2024ripple}. String-keyed editors get locality almost for free, whereas other editors need retain data~\citep{gangadhar2024finetuning}. \sdf propagates in natural language, egregious facts stay brittle, and \sdf causes the broadest drift in unrelated world knowledge~\citep{slocum2025believe,anthropic2025sdf}. Three findings are new for computed facts: the coverage effect of the regulariser, leakage that follows the answer format into counting, parity and sign handling, and the split between natural-language belief and equation computation under \sdf.

\para{Implications.} Locality benchmarks for editing and unlearning usually score a fixed set of neighbours. For any knowledge that the model computes rather than stores, such as arithmetic, unit conversion and possibly many rule-like facts, our results suggest that low neighbour change can co-exist with large changes elsewhere. Evaluations should include probes in other answer formats that no regulariser saw. Conversely, a perfect locality score on a broad suite should make us suspect that the edit is an exception, and should prompt a test of whether the edit survives a change of context.

\para{Limitations.} We study one base model (\qwen, digit-level tokens) with few-shot prompts. The 2+2$\to$5 target carries a strong cultural prior. We did not run MEMIT or AlphaEdit. Seeds mainly vary the LoRA initialization. Our \sdf corpus is small, and some control probe groups have only 4--16 items. The mechanistic analysis is correlational. \Appref{app:limitations} discusses each point in detail, along with broader impact.
